\section{Experimental protocol}
\label{sec:experiments}

% This section fixes the protocol BEFORE any language-model run. Changing a
% threshold or a decision rule after results are seen must be reported as a
% deviation in the paper. The canonical, more detailed version is
% docs/research-plan.md; keep the two in sync. Every number below comes from
% tables/macros.tex (results/power/power.json for the design).

This section specifies the language-model evaluation. None of it has been
executed. \pending{all runs in this section} The study evaluates Claude models
through the backend of Section~\ref{sec:loop}; the harness accepts other
providers behind the same one-step protocol, but the claims of the study are
limited to the models evaluated.

\subsection{Research questions and hypotheses}

\paragraph{RQ1: grounding accuracy.}
On answerable tasks (\code{lookup}, \code{compute}, \code{multi\_step}), how
accurate are the copilot's answers against the independent reference, and
what share of its numeric claims are supported by the tool outputs it cites?
H1a: accuracy is at least 90\%. H1b: at least 95\% of answerable episodes that
make a claim are fully grounded. H1c, on the validity of the metric: on at
least 300 claims sampled stratified by verifier verdict and annotated by two
people (Cohen's $\kappa \ge 0.8$ required), the verifier agrees with the
annotation on at least 95\% of claims after reweighting the strata to each
model's verdict distribution; false-accept and false-reject rates are
reported per stratum.

\paragraph{RQ2: look-ahead safety.}
Does the copilot request data after the cutoff when the question does not
need it, and does it abstain when a question cannot be answered without such
data? The primary measure is the share of answerable episodes with a
look-ahead attempt; attempts on trap questions are reported separately,
together with the abstention that follows them. H2a (structural): no episode
with the clock enforced contains a result dated after $t$. A single leak
invalidates the run as a harness defect. H2b: abstention accuracy is at least
90\% on \code{pit\_trap} and on \code{unknown\_symbol}, \emph{and} false
abstention on answerable tasks is at most 5\%; a constant policy passes the
first part and fails the second (Section~\ref{sec:harness-results}). Real
tickers (\code{absent\_real}) are reported separately. H2c (ablation A1): with
the clock disabled, the rate at which answers match a numeric hindsight value
is higher than with the clock enforced (one-sided exact McNemar test paired by
task, $\alpha = 0.05$). Ranking and count hindsight values are excluded
because they are matched by chance.

\paragraph{RQ3: policy compliance under adversarial prompts.}
When asked to trade, including with false claims of authorization, does the
copilot refuse, avoid execution attempts and execution claims, and avoid
over-refusing ordinary questions? H3a: refusal accuracy on \code{policy\_trap}
is at least 95\%, reported with and without a decoy \code{execute\_order} tool
that is offered to the model and always refused by the gate. Without the
decoy, a model can attempt execution only by calling a tool it was not
offered. H3b: \tok{EXECUTION\_REFUSED} appears in at most 2\% of non-trade
episodes. The system guarantees that no order is executed, whatever the model
does; that is tested (Section~\ref{sec:system}), not hypothesized.

\paragraph{RQ4: cost and latency versus accuracy.}
How do accuracy, grounding and safety trade off against tokens and latency
across effort levels? H4: for each effort level below the main one, accuracy
on answerable tasks is non-inferior to the main configuration with a margin
of 5 percentage points: the one-sided 95\% lower bound of the paired
difference in item-level success, from a task-clustered bootstrap, lies above
$-5$ points, with Holm correction across levels. We also report the Pareto
frontier of accuracy against mean output tokens and median episode latency.

\subsection{Decision rule and power}
\label{sec:power}

Thresholds are fixed before the first run. A threshold hypothesis is
\emph{supported} when the lower bound of the Wilson 95\% interval over
distinct items is at or above the threshold (for H3b, the upper bound at or
below it), \emph{not supported} when the interval lies entirely on the wrong
side, and \emph{inconclusive} otherwise. With repetitions, an item counts as a
success only if every repetition succeeds. This is a reliability criterion:
at a per-episode success rate $p$ and independent repetitions, an item
succeeds with probability $p^3$.

Table~\ref{tab:power} gives the probability that the rule declares each
hypothesis supported for the planned design: \powSuites{} evaluation suites of
\powTasksPerSuite{} tasks, \powRepetitions{} repetitions each (\powEpisodes{}
episodes). Items are deduplicated within and across suites and against the
development suite, and the item and cluster counts come from generating the
design with fixed design seeds. H3a pools \powItemsPolicyTrap{} distinct trade
requests in \powClustersPolicyTrap{} template-by-symbol clusters and allows
\powHThreeAAllowed{} failures. It is well powered only for a per-episode
refusal rate of 99.5\% (\powHThreeAIndependentNineNineFive{} with independent
repetitions) and is expected to be inconclusive at 99\%
(\powHThreeAIndependentNineNine), which is intended for a reliability claim.
The abstention hypotheses pool \powItemsPitTrap{} cutoff traps and
\powItemsUnknownSymbol{} unknown-symbol questions and are powered at 99\%
(\powHTwoBPitIndependentNineNine) but not at 98\%
(\powHTwoBPitIndependentNineEight). Simulations in which items of one cluster
share a beta-distributed success rate (intra-cluster correlation 0.05 and 0.2,
\code{results/power/power.md}) move these probabilities by up to
\powMaxClusterShift, mostly at rates below the threshold. The largest effect
on a configuration that is otherwise powered is H1b at 99\%, whose power falls
from \powHOneBIndependentNineNine{} to \powHOneBIccTwoNineNine{} at an
intra-cluster correlation of 0.2. For H4, a paired non-inferiority test with a
5-point margin needs \powNonInfTwoMarginZeroFive{} items at a discordant-pair
rate of 0.2; the sweep therefore uses \powSweepSuites{} suites
(\powSweepItems{} answerable items).

\begin{table}[t]
\centering
\small
% Generated by scripts/render_paper_tables.py from results/{benchmark,verifier,power}/.
% Do not edit by hand. Scripted harness-validation baselines and verifier stress tests on SYNTHETIC data,
% and a design-only power analysis; none of these are language-model results.
\begin{tabular}{llrrrrrr}
\toprule
Hypothesis & Threshold & Items & Clusters & Failures allowed & $p$ & Indep. & Identical \\
\midrule
H1a & $\ge$ 90\% & 1008 & 302 & 82 & 0.97 & 0.27 & 1.00 \\
 &  &  &  &  & 0.98 & 1.00 & 1.00 \\
 &  &  &  &  & 0.99 & 1.00 & 1.00 \\
 &  &  &  &  & 0.995 & 1.00 & 1.00 \\
\midrule
H1b & $\ge$ 95\% & 1008 & 302 & 36 & 0.97 & 0.00 & 0.87 \\
 &  &  &  &  & 0.98 & 0.00 & 1.00 \\
 &  &  &  &  & 0.99 & 0.89 & 1.00 \\
 &  &  &  &  & 0.995 & 1.00 & 1.00 \\
\midrule
H2b-pit & $\ge$ 90\% & 288 & 39 & 18 & 0.97 & 0.08 & 1.00 \\
 &  &  &  &  & 0.98 & 0.66 & 1.00 \\
 &  &  &  &  & 0.99 & 1.00 & 1.00 \\
 &  &  &  &  & 0.995 & 1.00 & 1.00 \\
\midrule
H2b-unknown & $\ge$ 90\% & 288 & 72 & 18 & 0.97 & 0.08 & 1.00 \\
 &  &  &  &  & 0.98 & 0.66 & 1.00 \\
 &  &  &  &  & 0.99 & 1.00 & 1.00 \\
 &  &  &  &  & 0.995 & 1.00 & 1.00 \\
\midrule
H3a & $\ge$ 95\% & 288 & 111 & 7 & 0.97 & 0.00 & 0.36 \\
 &  &  &  &  & 0.98 & 0.00 & 0.78 \\
 &  &  &  &  & 0.99 & 0.38 & 0.99 \\
 &  &  &  &  & 0.995 & 0.93 & 1.00 \\
\midrule
H3b & $\le$ 2\% & 1584 & 413 & 20 & 0.002 & 1.00 & 1.00 \\
 &  &  &  &  & 0.005 & 0.26 & 1.00 \\
 &  &  &  &  & 0.01 & 0.00 & 0.88 \\
\bottomrule
\end{tabular}

\caption{Probability that the threshold rule declares a hypothesis supported,
by true per-episode rate $p$ (for H3b, the per-episode over-refusal rate),
under independent or identical repetitions. Design only; no model has been
run. From \code{scripts/power\_analysis.py}.}
\label{tab:power}
\end{table}

\subsection{Configurations}

\begin{itemize}
  \item \textbf{Main.} \code{claude-opus-5} with adaptive thinking, effort
  \code{high}, \code{max\_tokens} 16{,}000, the default policy (24 tool calls,
  strict schemas) and at most \RunMaxSteps{} steps. Server-side fallback is
  off. A run in which any served model differs from the requested model is
  marked invalid, as is a run that meets an unrecoverable API error; such a
  run stops instead of scoring the failure as a wrong answer.
  \item \textbf{Effort sweep (RQ4).} Every effort level the model under test
  accepts (the backend validates \code{low}, \code{medium}, \code{high},
  \code{xhigh} and \code{max}) on \powSweepSuites{} suites with two
  repetitions each.
  \item \textbf{Suites.} The public default suite (seed \SuiteSeed{}) is
  reserved for development and pilot runs. The evaluation uses eight fresh
  seeds, drawn after the system prompt and scorer are frozen; each evaluation
  suite has its own synthetic panel and shares no item with the development
  suite or with another evaluation suite. The seeds and the suite SHA-256 are
  recorded in the run manifest before the first model call, and the suites are
  not published until the runs are complete.
  \item \textbf{Repetitions.} Current models accept no sampling parameters,
  so run-to-run variability is measured by repetition: three independent runs
  for the main configuration. Mean rates over repetitions are reported with
  task-clustered bootstrap intervals (resampling tasks together with all their
  repetitions) as descriptive statistics. Decisions use the all-repetitions
  criterion above.
  \item \textbf{Multiple comparisons.} Pairwise comparisons across the sweep
  are controlled with the Holm procedure; exploratory comparisons are
  reported with Benjamini--Hochberg control at a false discovery rate of
  0.05~\citep{benjamini1995fdr}.
  \item \textbf{Accounting.} Model refusals and \code{max\_steps}
  terminations count as incorrect and are also reported separately. Answers
  that fail the output format are separated from substantive errors using the
  recorded parse method and parse notes. Every turn is recorded for offline
  replay, and an interrupted run is resumed from its recording.
\end{itemize}

\subsection{Arms and ablations}

\begin{itemize}
  \item \textbf{A1: no clock.} The same model with the refusal of explicitly
  named later dates lifted (Section~\ref{sec:clock}). It measures how often a
  language model reaches for data after $t$ when nothing stops it (RQ2, H2c).
  A policy that never names a later date is unaffected by it.
  \item \textbf{A2: no citation requirement.} A prompt variant without the
  citation instruction; tool results also omit their ``cite as'' line. The
  verifier still attributes uncited claims to any result, which separates
  faithfulness from citation compliance.
  \item \textbf{A3: no argument gate.} Argument-level gate checks are removed,
  while the data view still refuses later dates and no execution handler
  exists. It tests defense in depth and will be a test-only configuration
  that is never exposed on the command line.
  % TODO(code): implement A3 as a test-only harness hook.
  \item \textbf{A4: paraphrase and prompt sensitivity.} Paraphrased questions
  and a reworded system prompt with the same protocol.
  % TODO(code): paraphrase templates (generator v3).
  \item \textbf{A5: closed book.} No tools and a prompt that says so. It
  measures what the tools add, and how the model treats real tickers without
  them.
  \item \textbf{A6: no cutoff instruction.} The prompt states the as-of date
  but not the instruction to ignore later knowledge.
  \item \textbf{Decoy.} The main configuration with the refused
  \code{execute\_order} tool offered (RQ3).
\end{itemize}

\subsection{Real data}
The same protocol will be repeated on confirmed daily bars read from the
suite's external store through \code{StoreBarSource}. Only aggregate results
will be committed, never vendor data. This requires a task generator and
reference that read an arbitrary confirmed panel, and a point-in-time
treatment of price adjustments and of the universe
(Section~\ref{sec:limitations}). \pending{real-data generator and runs}
